\documentclass[10pt]{amsart}

\usepackage[T1]{fontenc}
\usepackage{lmodern}
\usepackage{amsmath,amssymb,amsthm,mathtools}
\usepackage[margin=1.08in]{geometry}
\usepackage[final]{microtype}
\usepackage[colorlinks=true,linkcolor=blue,citecolor=blue,urlcolor=blue]{hyperref}

\newtheorem{theorem}{Theorem}[section]
\newtheorem{proposition}[theorem]{Proposition}
\newtheorem{corollary}[theorem]{Corollary}
\theoremstyle{definition}
\newtheorem{definition}[theorem]{Definition}
\theoremstyle{remark}
\newtheorem{remark}[theorem]{Remark}

\newcommand{\C}{\mathbb C}
\newcommand{\U}{\mathcal U}
\newcommand{\Shadow}{\mathfrak S}
\newcommand{\MFRad}{\operatorname{Rad}_{\mathrm{MF}}}
\newcommand{\Normal}[1]{\langle\!\langle #1\rangle\!\rangle}

\hypersetup{
  pdftitle={Shadow--Kleene Saturation},
  pdfauthor={Sauers},
  pdfsubject={From one tracial bug to a full MF radical},
  pdfkeywords={MF groups, property (T), Hilbert--Schmidt ultraproducts,
    acylindrical small cancellation, Kleene recursion theorem}
}

\title[Shadow--Kleene saturation]{Shadow--Kleene Saturation:\\
  From One Tracial Bug to a Full MF Radical}
\author{Sauers}
\date{August 22, 2026}

\begin{document}

\begin{abstract}
We isolate a weak intermediate obstruction to operator MF: a nonidentity
word that vanishes in the normalized Hilbert--Schmidt shadow of every
operator-norm almost representation.  Higman embedding and acylindrical
small cancellation compile the normal closure of one such word onto a
2-generated finitely presented Kazhdan group.  A normal-Kazhdan moving-corner
detector then upgrades shadow invisibility to full operator-MF invisibility.
The output has full MF radical, as does every quotient.  A double-HNN
compression supplies an unconditional torsion-free seed, while a
threshold-free reverse-Kleene theorem reduces the computational program to
constructing one shadow word.
\end{abstract}

\maketitle

\section{The shadow residual}

For dimensions $d_n$, put
\[
 \mathcal Q_{\mathrm{op}}
 =\prod_nM_{d_n}(\C)\big/
   \bigoplus_n^{\|\cdot\|_{\mathrm{op}}}M_{d_n}(\C).
\]
The operator-MF radical of a countable group $G$ is
\[
 \MFRad(G)=\bigcap_\Theta\ker\Theta,
\]
where $\Theta$ ranges over homomorphisms from $G$ to unitary groups of norm
matrix coronas.  Thus $\MFRad(G)=G$ precisely when every such homomorphism is
trivial.

\begin{definition}
The \emph{operator-to-HS shadow residual} $\Shadow(G)=R_{\infty\to2}(G)$
consists of the $w\in G$ such that, for every sequence
$\phi_n:G\to\U(d_n)$ with
\[
 \|\phi_n(xy)-\phi_n(x)\phi_n(y)\|_{\mathrm{op}}\longrightarrow0
 \qquad(x,y\in G),
\]
the image of $w$ is $1$ in every normalized Hilbert--Schmidt ultraproduct.
A nonidentity element of $\Shadow(G)$ is a \emph{shadow bug}.
\end{definition}

Since normalized HS norm is bounded by operator norm,
$\MFRad(G)\leq\Shadow(G)$.  This inclusion can be strict even for an MF
group; a shadow bug is therefore not a disguised non-MF hypothesis.

\begin{proposition}[Functoriality]
$\Shadow(G)$ is a fully invariant normal subgroup.  For every homomorphism
$f:G\to H$,
\[
 f(\Shadow(G))\leq\Shadow(H).
\]
Consequently $w\in\Shadow(G)$ implies
$\Normal{w}^{G}\leq\Shadow(G)$.
\end{proposition}

\begin{proof}
An operator-norm almost representation of $H$ restricts along $f$ to one of
$G$.  Normality follows either by applying this observation to inner
automorphisms or by writing $\Shadow(G)$ as an intersection of kernels.
\end{proof}

\section{The moving-corner detector}

\begin{theorem}[Normal-Kazhdan detector]\label{thm:detector}
Let $K\trianglelefteq G$ have property~\textup{(T)}.  If
$K\leq\Shadow(G)$, then
\[
 K\leq\MFRad(G).
\]
In particular, if $G$ has property~\textup{(T)} and $\Shadow(G)=G$, then
$\MFRad(G)=G$.
\end{theorem}

\begin{proof}[Mechanism]
Suppose a norm-corona representation keeps an element of $K$ alive.  Its
coordinate lifts may move only a subspace of vanishing normalized rank, so
the original HS shadow need not see it.  A Kazhdan Laplacian cuts out the
moving projection.  Normality of $K$ makes these projections asymptotically
invariant under the ambient group.  Compression and polar correction give
another operator-norm almost representation, now normalized by the moving
rank.  The Kazhdan inequality forces one element of a fixed Kazhdan set to
survive in its HS ultraproduct, contradicting $K\leq\Shadow(G)$.
\end{proof}

This detector replaces the finite spectral corner of a central involution
by the moving-space projection of an infinite Kazhdan subgroup.  It uses no
central sign, Clifford representation, or Hilbert-hotel construction.

\section{Saturating one word}

\begin{theorem}[One-word shadow saturation]\label{thm:saturation}
Let $G$ be finitely generated and recursively presented.  If
\[
 1\neq w\in\Shadow(G),
\]
then there is a 2-generated finitely presented acylindrically hyperbolic
property-\textup{(T)} group $Q$ such that
\[
 \MFRad(Q)=Q.
\]
If the source is already finitely presented, torsion-free, and
acylindrically hyperbolic, $Q$ may be chosen torsion-free.
\end{theorem}

\begin{proof}
Effective Higman embedding~\cite{Higman} places $G$ in a finitely presented group $H$;
functoriality preserves both the shadow property and, by injectivity,
nontriviality of $w$.  Set $E=H*F_2$ and
$N=\Normal{w}^{E}$.  The Bass--Serre action makes $E$ non-elementary
acylindrically hyperbolic with trivial finite radical.  The subgroup $N$ is
infinite, normal, and hence a suitable reservoir by the acylindrical subgroup
theory of Osin~\cite{Osin}.

Take a fixed finitely presented non-elementary hyperbolic Kazhdan group
$H_T$.  Run Hull's relative common-quotient construction with two
noncommensurable loxodromics $h_1,h_2\in N$~\cite{Hull}, protect $w$, and route every
generator $x_i$ of $E$ by a long relation
\[
 x_i=W_i(h_1,h_2).
\]
The resulting common quotient $E\twoheadrightarrow Q\leftarrow H_T$ is
2-generated, finitely presented, acylindrically hyperbolic, Kazhdan, and
satisfies
\[
 q(w)\neq1,
 \qquad q(N)=Q.
\]
Hull's torsion clause gives the final assertion in the torsion-free case.

Because $N\leq\Shadow(E)$, functoriality yields
$Q=q(N)\leq\Shadow(Q)$.  Hence $\Shadow(Q)=Q$, and
Theorem~\ref{thm:detector} gives $\MFRad(Q)=Q$.
\end{proof}

\begin{corollary}\label{cor:quotients}
Every quotient of $Q$ has full MF radical.  Consequently every nontrivial
quotient is non-MF and every homomorphism from $Q$ to an MF group is trivial.
\end{corollary}

\begin{proof}
A norm-corona representation of a quotient pulls back to one of $Q$.  For a
map $Q\to M$ with $M$ MF, its image is both a quotient of $Q$ and an MF
subgroup, hence trivial.
\end{proof}

If $Q=\langle a,b\mid R\rangle$ and $v\neq1$ in $Q$, the finite conditions
$R=1$ and $v\neq1$ define a nonempty clopen family of 2-marked non-MF groups:
each member is a nontrivial quotient of $Q$.  This is a finite challenge
certificate for full MF failure.

\section{An unconditional torsion-free seed}

Let $P$ be finitely presented, torsion-free, and Kazhdan, containing
\[
 P_1\times P_2\times S\leq P,
\]
where $P_i\cong P$ and $S$ is infinite, finitely presented, nonabelian,
simple, and torsion-free.  Form the double HNN extension
\[
 E=\langle P,u_1,u_2\mid u_iPu_i^{-1}=P_i,\ i=1,2\rangle.
\]
The Fournier--Facio skeleton~\cite{FF} gives $E$ finitely presented,
torsion-free, and acylindrically hyperbolic.  Choose $s,x\in S$ with
$[s,x]\neq1$ and put
\[
 t=u_1,\qquad c=u_1^{-1}su_1,\qquad w=[s,x].
\]
Then $[c,P]=1$ and $tct^{-1}=s$.  Britton's lemma embeds $P$, so $w\neq1$
in $E$.  The Kazhdan compression theorem transports the asymptotic
centralizer of $P_1=tPt^{-1}$ to that of $P$ in every normalized-HS shadow
of an operator-norm almost representation.  Hence $s$ commutes there with
$x\in P$, and
\[
 1\neq w\in\Shadow(E).
\]
Applying Theorem~\ref{thm:saturation} directly yields a 2-generated,
finitely presented, torsion-free, acylindrically hyperbolic Kazhdan group
$Q$ with $\MFRad(Q)=Q$.

\section{The Shadow--Kleene front end}

\begin{theorem}[Shadow--Kleene]
Suppose a total computable compiler sends a machine index $e$ to a uniformly
recursively presented group $\Gamma_e$ and a word $w_e$, with
\[
 M_e\downarrow\ \Longrightarrow\ w_e\neq1,
 \qquad
 M_e\uparrow\ \Longrightarrow\ w_e\in\Shadow(\Gamma_e).
\]
Then a 2-generated finitely presented acylindrically hyperbolic Kazhdan
group with full MF radical exists.
\end{theorem}

\begin{proof}
By Kleene's recursion theorem, construct a machine which computes its own
compiled pair and halts exactly when proof enumeration finds $w=1$.  It
cannot halt, since the displayed HALT clause would contradict the proof it
found.  Nonhalting also implies $w\neq1$: equality in a recursively
presented group has a finite derivation and is recursively enumerable.  The
NONHALT clause now gives a nontrivial shadow bug, to which
Theorem~\ref{thm:saturation} applies.
\end{proof}

The fixed-point machine searches proofs only.  It reads no matrices,
dimensions, approximation thresholds, convergence rates, or soundness
moduli.  The remaining computational target is therefore the qualitative
implication
\[
 \text{NONHALT}\Longrightarrow w_e\in R_{\infty\to2}(\Gamma_e),
\]
which is weaker than erasure in all HS microstates.

\section{Formalization boundary}

The Lean file
\path{GroupApproximation/Sofic/OpToHSShadowResidual.lean} defines
$\Shadow(G)$ and checks its normality and functoriality, the passage from one
word to its normal closure, the normal-Kazhdan detector interface, and the
analytic saturation endpoint
\path{normMFResidual_eq_top_of_shadow_bug_saturation}.  It also
extracts the explicit HNN commutator via
\path{FournierFacioDefectData.exists_nontrivial_opToHSShadowBug}.
Effective Higman embedding, Bass--Serre/Osin suitability, and Hull's
protected common-quotient construction are the stated external
group-theoretic boundary.  The compression-free computational seed remains
open.

\begin{thebibliography}{9}

\bibitem{FF}
F.~Fournier--Facio,
\emph{A torsion-free non-sofic group},
\href{https://arxiv.org/abs/2608.02025}{arXiv:2608.02025} (2026).

\bibitem{Higman}
G.~Higman,
\emph{Subgroups of finitely presented groups},
Proc. Roy. Soc. Ser. A 262 (1961), 455--475.

\bibitem{Hull}
M.~Hull,
\emph{Small cancellation in acylindrically hyperbolic groups},
Groups Geom. Dyn. 10 (2016), 1077--1119.

\bibitem{Osin}
D.~Osin,
\emph{Acylindrically hyperbolic groups},
Trans. Amer. Math. Soc. 368 (2016), 851--888.

\end{thebibliography}

\end{document}
